% THIS IS A LATEX TEMPLATE FILE FOR PAPERS INCLUDED IN THE
% *Anthology of Computers and the Humanities*.
\documentclass[final]{anthology-ch}

% LOAD LaTeX PACKAGES
\usepackage{booktabs}
\usepackage{graphicx}

% TITLE OF THE SUBMISSION
\title{Opening Digital Humanities to All: Approaches to Supporting Accessible Omeka Projects}

% AUTHOR AND AFFILIATION INFORMATION
\author[1]{Allison Kaefring}[
  orcid=0009-0006-3805-7303
]

\author[1]{Jennifer Daugherty}[
  orcid=0009-0001-3189-7944
]

\author[1]{Jeanne Hoover}[
  orcid=0000-0003-3259-3640
]

\affiliation{1}{East Carolina University, Greenville, NC, USA}

% KEYWORDS
\keywords[eng]{accessibility, digital humanities, Omeka, WCAG 2.1, ADA Title II, web accessibility}

% METADATA FOR THE PUBLICATION
\pubyear{2025}
\pubvolume{5}
\pagestart{58}
\pageend{62}
\paperorder{7}
\conferencename{From Manuscript to Megabyte: Building Sustainable and Open Futures in the Digital Humanities}
\conferenceeditors{James B. Harr III, Emma Stanley, Jenna Rinalducci, and Donna Kain}
\doi{10.63744/PkXgYsplgWMW}
\customhead{}

\addbibresource{bibliography.bib}

%%%%%%%%%%%%%%%%%%%%%%%%%%%%%%%%%%%%%%%%%%%%%%%%%%%%%%%%%%%%%%%%%%%%%%%%%%%
% HERE IS THE START OF THE TEXT
\begin{document}

\maketitle

\begin{abstract}
In 2024, the Department of Justice issued revised Title II regulations under the American Disabilities Act, requiring that all government institutions make their digital content accessible and in compliance with WCAG 2.1 AA guidelines. In response, the Digital Humanities Working Group at a public institution is navigating the procedural changes needed to ensure our future digital humanities projects are fully accessible, while also exploring technical and ethical ways to make our past projects accessible. Using tools like WAVE, we have evaluated common accessibility barriers to our past digital humanities projects based in Omeka. The working group is developing solutions to our future digital humanities tool trainings, workflows, and user agreements. In this paper, we will share our process, the challenges encountered, and recommendations or supporting students and faculty in creating digital humanities projects that are accessible to all users. This paper will address both common accessibility issues we have identified in digital humanities projects and ways to make our commitment to accessibility sustainable.
\end{abstract}

\section{Introduction} \label{sec:intro}

Since the new American with Disabilities Act (ADA) Title II regulations \cite{doj2024ada} were published in 2024, East Carolina University (ECU) and ECU Libraries have been systematically working through websites and digital resources to make them accessible. Recently, the library's Digital Humanities Working Group (DHWG) explored ways to make digital humanities projects accessible, particularly Omeka S projects. Additionally, the DHWG reviewed current Omeka S projects on accessibility and identified common issues.

\section{Accessibility and Digital Humanities} \label{sec:accessibility}

New \href{https://www.ada.gov/law-and-regs/regulations/title-ii-2010-regulations/}{Americans with Disabilities Act (ADA) Title II regulations} \cite{doj2024ada} published in June 2024 require all government entities to make their web content accessible to \href{https://www.w3.org/TR/WCAG21/}{WCAG (Web Content Accessibility Guidelines) 2.1 AA} \cite{w3c2025wcag} standards. These guidelines were created by W3C (World Wide Web Consortium), an organization dedicated to web accessibility. Originally, all content needed to be in compliance with WCAG 2.1 AA by April 2026, however, the deadline was extended to April 2027.

The WCAG standards are composed of 4 main sections: Perceivable, operable, understandable, and robust \cite{w3c2025wcag}. Operable, understandable, and robust accessibility components primarily relate to software functionality and a user's ability to navigate the site. These include aspects of a website's code that allows for keyboard navigation, proper time to navigate a site without it timing out, predicable outcomes for actions, interoperability with assistive technology, and input assistance to help users avoid or correct mistakes. As such, most of these components are outside of the scope of the average digital humanities project that uses existing software. So long as digital humanities projects use software with a VPAT (Voluntary Product Accessibility Template) or ACR (Accessibility Compliance Report) that states that the software is compliant with WCAG 2.1 AA standards, no further intervention is needed for the vast majority of these components.

The few notable areas where the other components do affect digital humanities projects that utilize existing software include page navigation order, explanations for abbreviations, and avoiding screen flashes that could induce seizures. These components could become more relevant if a digital humanities scholar chose to edit the underlying code of the software. While the rise in AI \enquote{vibe coding,} where artificial intelligence writes the code instead of a human, has made it increasingly easy for people with minimal coding experience to make changes to the underlying code of existing software to customize it to their project \cite{williams2025vibe}. However, the lack of programming knowledge runs the risk of removing or fundamentally changing aspects of the code that were created to adhere to these guidelines. In cases where Digital Humanities scholars are interested in altering the code of a software, either by hand or through vibe coding, it is important to ensure that the end result still meets WCAG 2.1 AA guidelines so researchers or students who use assistive technology are still able to engage with the project.

Perceivable components, however, relate to properly using the software's built-in accessibility features to ensure the content is accessible. Perceivable components include things like alternative text for images, captions for videos, appropriate color contrast and size for text, and properly nested headings. These components play a vital role in ensuring that people using assistive technologies are able to gain the same amount of information from a Digital Humanities project as other users.

At ECU all of the library's website updates to ensure digital accessibility were handled by the library's Information Technology Department (IT). However, there were a few gaps in their accessibility assessment for web content created by the library that was technically not hosted on the library webpages. This included digital library collections like the institutional repository, LibGuides, and our digital humanities projects in hosted software such as Omeka.

The only exceptions for ensuring web content comply with WCAG 2.1 AA standards are for archived content, personalized and password protected content, or preexisting content not currently in use \cite{doj2024ada}. This presents an interesting quandary for digital humanities projects that are already completed. It is difficult to assess whether or not a project such as a completed student capstone project that lives in an Omeka site would be considered \enquote{archived.} Many of these sites were built with the intention of creating resources that could continue to be used by the community---which would imply that are (ideally) still currently in use. However, for most of the Omeka capstones and class projects, the creators of the sites themselves graduate fairly soon after the project is complete. The library could consider these Omeka sites as \enquote{archived} since they are no longer being updated, and often, their student creators no longer have access to make updates.

\section{Digital Humanities and Omeka Projects at ECU} \label{sec:omeka}

Early digital humanities work at ECU included the development of the Digital Innovation and Scholarship in Social Sciences and Humanities (DISSH) initiative \cite{ecudissh}. Driven by faculty, the initiative strove to bring digital humanities concepts and tools to the broader campus. During that time, the Information Technology and Computing Services Division at ECU utilized Omeka Classic managed hosting of version 3.1.1. Due to changes at the university's administrative level and faculty attrition, DISSH lost momentum and support.

In 2021, ECU Libraries' internal IT team took over management of Omeka Classic and switched to self-hosting the Omeka S platform. There was an attempt to migrate a limited number of projects over to the library platform if they were in active use. Some of these projects are older than ten years and would be difficult to track down the students or the faculty who originated them. The rest of the projects remain on a server and are visible, but there is very limited ability to make edits. Due to their status, that adds weight to the argument that they could be considered archived sites under the web accessibility guidelines. The library's Digital Humanities Working Group (DHWG) now manages all requests for an Omeka site. The DHWG consists of a cross-departmental library team that provides support to faculty and staff for projects. They also host lunch and learn workshops aimed at encouraging digital humanities engagement. Faculty and others can submit a form to request a site to the DHWG. The working group provides support throughout the process including providing instruction on use of Omeka to the requester or their students and instruction on copyright, troubleshooting issues that might arise, and helping make decisions regarding sustainability of projects and archiving projects when needed or requested. The University Archives also uses the Archive-It tool to archive the front facing aspects of the sites.

The bulk of the current Omeka requests come from the History Department. The department has approved a digital humanities capstone project to substitute for the required undergraduate senior thesis or for the graduate thesis, which has encouraged use of the Omeka platform. A recent graduate student did a project, \enquote{South Atlantic Powwow Archive,} focusing on the evolution of public pow-wows in Native American culture. Available at \url{https://collections.ecu.edu/os/s/SouthAtlanticPowwowArchive/page/About-It}. Another graduate student did a project, \enquote{Black History of Ayden,} focusing on the Black history of a small rural community in Pitt County, NC. Available at \url{https://collections.ecu.edu/os/s/black-history-of-ayden/page/home}. These exemplify the typical structure of a project.

There has also been an increase in faculty and instructors using OmekaS for class-wide projects. The instructors in the new \textit{Archives and Cultural Heritage Informatics} focus in the graduate library science program at ECU have incorporated Omeka usage into several of their classes and have communicated a desire to carry over a sandbox site from one semester to another. The transitory use of the platform for pedagogical instruction might create additional technical and accessibility issues.

Other digital humanities tools are available but are not as widely used. Tools like ArcGIS's StoryMaps are much more complex and can make it difficult for their integration into a semester long course or for students to learn on their own for capstone or thesis projects. Omeka can be quickly taught, and students can use it almost immediately. The DHWG can also make many requested changes to the platform without IT's involvement. When IT must be involved, it is usually a quick process since they understand the product, and it is hosted by the library.

\section{Current and Future Plans} \label{sec:plans}

With the new accessibility regulations, the Digital Humanities Working Group discussed how to best make current and future Omeka projects accessible. The first step was to assess the accessibility of current Omeka projects. The DHWG used a free tool, WAVE, developed by WebAIM at Utah State University \cite{webaimwave} to assess the accessibility of current Omeka projects. The DHWG decided to focus on the \enquote{errors} and \enquote{alerts} that were noted in the WAVE tool's accessibility summary. The tool provides additional information in the accessibility summary, but the common accessibility issues were captured by the errors and alerts section. Most of the current Omeka projects only had a handful of errors or alerts. For example, one OmekaS site had an issue with color contrast and needed alternative text for the images. The most common errors across the OmekaS sites were related text size, color contrast, missing alternative text for images, missing closed captions or transcriptions, nested headings, and properly formatted lists. The review of current Omeka sites informed what next steps the Digital Humanities Working Group needed to do to support the accessibility of digital humanities projects going forward.

The DHWG's top priority is developing guidance on making accessible Omeka projects. This includes incorporating information into library instruction sessions and developing instructions and videos for the library's digital humanities guide. Additionally, a simple Omeka accessibility checklist will also be provided to future project creators. These resources will also be combined with conversations with creators about the importance of accessible projects.

For the past few years, faculty involved with digital humanities along with the Digital Humanities Working Group have hosted a lunch and learn series during the academic year. The series provides opportunities for faculty and students to share their research projects or share how they are doing digital humanities in their courses. Sessions also cover various digital humanities tools and methodologies. At least one session will discuss accessibility at a future lunch and learn during the upcoming year.

Finally, the Digital Humanities Working Group will need to adjust the library's Omeka procedures to include accessibility. On a yearly basis, the Digital Humanities Working Group will start reviewing Omeka projects for accessibility using the WAVE tool. The review will be time consuming, and the projects need to be divided amongst DHWG members. Currently, the free WAVE tool will only work on each page within Omeka. Most of the current projects at ECU contain multiple pages per project. While this process is not ideal, it will help the DHWG evaluate the accessibility of projects. The DHWG will continue to review accessibility tools to determine if a different tool would work better than WAVE for these projects.

In addition to reviewing the OmekaS projects, the Digital Humanities Working Group will revise their Memorandum of Understanding (MOU) for Omeka. The current Omeka MOU does not include information on accessibility of projects. Like sustainability of projects, it is important to discuss accessibility with faculty and students using the library's Omeka software. Incorporating a new section on accessibility within the MOU will provide the ability for Omeka users and librarians to develop an agreement regarding accessibility. Additionally, the revised agreement will include wording to allow the library to migrate or adapt a project, without changing content, for accessibility. The library's institutional repository license agreement has similar wording, although not specific to accessibility. The license notes, \enquote{You grant to ECU the right, without changing the content, to migrate one or more copies of the submission to any medium or format for security and preservation purposes.} \cite{ecuscholarship}. The Digital Humanities Working Group will use a variation of this sentence to include accessibility, as well as security, for access purposes.

While the new accessibility regulations have been delayed, the working group is committed to addressing accessibility of Omeka sites going forward. The DHWG wants to support faculty and students in making their digital humanities projects as accessible as possible. This can be achieved by revising our instructions and procedures for digital humanities support in the library.

\section{Considerations} \label{sec:considerations}

With the new guidelines and restrictions that must be legally adhered to, there comes more risk. With the onus being on the professor to make sure the project meets standards, they take on added responsibility. There are concerns that this may dissuade future creation of Omeka and other digital humanities projects. By offering training and guidance, the DHWG hopes to mitigate some of these concerns.

There may also be concerns from students regarding restrictions on their own creative activities. The DHWG plans to encourage the use of accessible templates for the OmekaS sites, and students may not agree with the color scheme or fonts. As mentioned before, even students with no technical background have learned to us AI to write code to modify the sites to make them more visually appealing, which then may render them non-compliant. This could provide a valuable opportunity for the DHWG to collaborate with instructors in delivering classroom training on the accessibility standards for OmekaS as part of course instruction. This would allow students to continue to modify their sites but also make them aware of the necessity and reasoning behind certain restrictions. It might also discourage the modification of coding changes that might affect underlying accessibility structures.

\section{Conclusion} \label{sec:conclusion}

One of the benefits of digital humanities projects is that they can be accessed by researchers and students all over the world. The reach of these projects can be increased by applying accessibility rules and best practices. Omeka is one of the most popular tools for digital humanities projects at ECU. The library's Digital Humanities Working Group supports Omeka along with other library-related digital humanities projects. The group will also provide the resources needed that will allow instructors and students to continue usage while also following necessary federal regulations. The result will not only maintain the university's compliance, but it will also widen the audience of those who can interact with the projects and engage with their content.

\printbibliography

\end{document}
